\section{体积弹性模量的估算与应用}

\begin{tipbox}{关键词标签}
体积弹性模量、等温压缩、压强估算、小量近似。
\end{tipbox}

\subsection{估计压缩到正常密度 $\alpha$ 倍所需的压强}

\textbf{题目}：估计将一固体压缩到它的正常密度 $\alpha$ 倍时所需的压强。

\textbf{分析}：本题是体积弹性模量 $K$ 的\textbf{定义式直接应用}，关键在于把微分定义改写为可积分的形式，并合理近似。

\vspace{0.5em}
\textbf{【解答】}

体积弹性模量的定义为
\[
  K = -V\left(\frac{\partial p}{\partial V}\right)_T .
\tag{1}
\]

在等温过程中把偏导改写为微分：
\[
  \dd p = -K\,\frac{\dd V}{V} .
\tag{2}
\]

假设 $K$ 在压缩过程中近似为常数（小形变近似），两边积分：
\[
  \int_{p_0}^{p}\dd p = -K\int_{V_0}^{V}\frac{\dd V}{V}
  \quad\Longrightarrow\quad
  p - p_0 = -K\ln\frac{V}{V_0} .
\tag{3}
\]

固体在自然平衡状态下仅受大气压作用，$p_0 \sim 10^5\,\text{Pa}$，而固体 $K \sim 10^{10}\,\text{Pa}$，故 $p_0 \ll K$ 可忽略，取 $p_0 \approx 0$：
\[
  p = -K\ln\frac{V}{V_0} .
\tag{4}
\]

密度变为正常密度的 $\alpha$ 倍，即 $\rho = \alpha\rho_0$。质量不变时 $\rho \propto 1/V$，故
\[
  \frac{V}{V_0} = \frac{1}{\alpha} .
\tag{5}
\]

代入 (4) 得
\[
  \boxed{p = K\ln\alpha} .
\tag{6}
\]

\begin{insightbox}{数量级估算}
典型离子晶体 $K \sim 10^{10}\,\text{Pa}$。若 $\alpha = 10$（压缩到 10 倍正常密度），则
\[
  p = 10^{10}\ln 10 \approx 2.3\times 10^{10}\,\text{Pa} \sim 10^5\,\text{atm} .
\]
这相当于地幔深度的压强，说明要把固体压缩一个数量级需要\textbf{天文数字的压强}。
\end{insightbox}

\begin{modelbox}{公式 (6) 的物理含义}
压强与 $\ln\alpha$ 成正比，而不是与 $(\alpha-1)$ 成正比——这说明体积弹性模量的定义本身就包含了\textbf{对数响应}。如果误用胡克定律 $p = K(1-V/V_0) = K(1-1/\alpha)$，会在 $\alpha \gg 1$ 时严重低估所需压强。

正确记忆：\textbf{$K$ 来自对数积分，结果自带 $\ln$}。
\end{modelbox}


\subsection{NaCl 在高压下离子间距的缩小百分比}

\textbf{题目}：已知 NaCl 晶体的体弹性模量为 $K = 2.4\times 10^{10}\,\text{Pa}$，在 $2$ 万大气压的压力作用下，晶体中相邻两离子间的距离将缩小百分之几？（$1\,\text{atm} = 10^5\,\text{Pa}$）

\textbf{分析}：上一题给出了 $p$ 与 $V$ 的关系，而离子晶体中 $V \propto r^3$（各向同性压缩），因此可以建立 $p$ 与 $r/r_0$ 的关系，再利用 $p \ll K$ 做泰勒展开。

\vspace{0.5em}
\textbf{【解答】}

由上题结论 (4)：
\[
  p = -K\ln\frac{V}{V_0} .
\tag{7}
\]

各向同性压缩时，晶胞体积与离子间距的立方成正比：
\[
  \frac{V}{V_0} = \left(\frac{r}{r_0}\right)^{\!3} .
\tag{8}
\]

代入 (7)：
\[
  p = -K\ln\!\left(\frac{r}{r_0}\right)^{\!3}
    = -3K\ln\frac{r}{r_0} .
\tag{9}
\]

解出间距比：
\[
  \ln\frac{r}{r_0} = -\frac{p}{3K}
  \quad\Longrightarrow\quad
  \frac{r}{r_0} = \exp\!\left(-\frac{p}{3K}\right) .
\tag{10}
\]

间距缩小的百分比为
\[
  \eta = \frac{r_0 - r}{r_0} = 1 - \frac{r}{r_0}
       = 1 - \exp\!\left(-\frac{p}{3K}\right) .
\tag{11}
\]

\textbf{代入数据}：
\[
  p = 2\times 10^4\,\text{atm}
    = 2\times 10^4 \times 10^5\,\text{Pa}
    = 2\times 10^9\,\text{Pa} .
\]

\[
  \frac{p}{3K} = \frac{2\times 10^9}{3\times 2.4\times 10^{10}}
               = \frac{2}{72}
               = \frac{1}{36}
               \approx 2.778\times 10^{-2} .
\tag{12}
\]

由于 $p \ll 3K$，对指数做泰勒展开 $e^{-x} \approx 1 - x + x^2/2$（$x = 1/36 \ll 1$）：
\[
  \frac{r}{r_0} = e^{-1/36}
               \approx 1 - \frac{1}{36} + \frac{1}{2\times 36^2}
               = 1 - 0.02778 + 0.000386
               \approx 0.9726 .
\]

因此
\[
  \boxed{\eta = 1 - 0.9726 \approx 2.74\%} .
\]

若只取一阶近似 $\eta \approx p/(3K) = 1/36 \approx 2.78\%$，误差仅约 $0.04\%$。

\begin{insightbox}{热力学符号与变量依赖关系}

热力学中符号混乱的根源是\textbf{两类物理量}的混淆：
\begin{itemize}[nosep]
  \item \textbf{状态函数}（如 $p,V,T,U$）：只与当前状态有关，与怎么到达无关；
  \item \textbf{过程量}（如 $W,Q$）：与具体路径有关，不能只说初末态。
\end{itemize}

\textbf{符号规则}：
\begin{center}
\renewcommand{\arraystretch}{1.4}
\begin{tabular}{|l|l|l|}
\hline
\textbf{符号} & \textbf{含义} & \textbf{例子} \\
\hline
$\partial$（偏导） & 固定其他变量求导 & $(\partial p/\partial V)_T$ 表示"温度不变时压强随体积的变化率" \\
\hline
\mbox{\textup{d}}（全微分） & 状态函数的无穷小变化 & $\mathrm{d}U$ 是内能变化，$\oint\mathrm{d}U = 0$ \\
\hline
$\delta$（变分） & 过程量的无穷小量 & $\delta W = p\,\mathrm{d}V$（功与路径有关，$\oint\delta W \neq 0$） \\
\hline
\end{tabular}
\end{center}

\vspace{0.5em}
\textbf{具体到本题}：$K = -V(\partial p/\partial V)_T$ 中的偏导数，为什么能直接积分？

因为 $p$ 是\textbf{状态函数}，在等温条件下 $p$ 只是 $V$ 的函数 $p = p(V)$（$T$ 固定）。此时偏导数退化为普通导数：
\[
  \left(\frac{\partial p}{\partial V}\right)_T = \frac{\mathrm{d}p}{\mathrm{d}V} \quad (T\text{ 固定}) .
\]

于是 $\mathrm{d}p = (\partial p/\partial V)_T\,\mathrm{d}V = -K\,\mathrm{d}V/V$，两边积分即可。

\textbf{什么时候会出错}：如果题目说"绝热压缩"，那么 $p$ 同时依赖于 $V$ 和 $S$（或 $T$），$p = p(V,S)$。此时
\[
  \mathrm{d}p = \left(\frac{\partial p}{\partial V}\right)_S\mathrm{d}V + \left(\frac{\partial p}{\partial S}\right)_V\mathrm{d}S ,
\]
若盲目只写 $\mathrm{d}p = (\partial p/\partial V)_S\,\mathrm{d}V$ 而忘了 $\mathrm{d}S \neq 0$，就会漏项。

\vspace{0.5em}
\textbf{热力学量为什么关注熵、热容、化学势}：

\begin{itemize}[nosep]
  \item \textbf{熵 $S$}：衡量系统\textbf{混乱程度}，也是热量传递的"方向标"——$\mathrm{d}Q_{\text{rev}} = T\,\mathrm{d}S$。绝热过程中 $S$ 不变，等温过程中 $S$ 随吸热而增加；
  \item \textbf{热容 $C$}：衡量系统\textbf{储热能力}，即温度升高一度需要多少能量。$C_V$ 和 $C_p$ 的区别在于约束条件（定容 vs 定压）；
  \item \textbf{化学势 $\mu$}：粒子进出的\textbf{"驱动力"}。平衡时各相的化学势相等（$\mu_1 = \mu_2$），就像热平衡时温度相等（$T_1 = T_2$）。$\mu$ 高的一侧粒子会向 $\mu$ 低的一侧流动。
\end{itemize}

\vspace{0.5em}
\textbf{自然变量与依赖关系}：

热力学中每个量该当作谁的函数，由\textbf{自然变量}决定。\textbf{关键原则}：同一个量在不同约束条件下有不同的自然变量。

\begin{center}
\renewcommand{\arraystretch}{1.4}
\begin{tabular}{|c|c|c|}
\hline
\textbf{热力学势} & \textbf{自然变量} & \textbf{基本关系} \\
\hline
内能 $U$ & $S, V, N$ & $\mathrm{d}U = T\,\mathrm{d}S - p\,\mathrm{d}V + \mu\,\mathrm{d}N$ \\
焓 $H$ & $S, p, N$ & $\mathrm{d}H = T\,\mathrm{d}S + V\,\mathrm{d}p + \mu\,\mathrm{d}N$ \\
自由能 $F$ & $T, V, N$ & $\mathrm{d}F = -S\,\mathrm{d}T - p\,\mathrm{d}V + \mu\,\mathrm{d}N$ \\
吉布斯函数 $G$ & $T, p, N$ & $\mathrm{d}G = -S\,\mathrm{d}T + V\,\mathrm{d}p + \mu\,\mathrm{d}N$ \\
巨势 $\Omega$ & $T, V, \mu$ & $\mathrm{d}\Omega = -S\,\mathrm{d}T - p\,\mathrm{d}V - N\,\mathrm{d}\mu$ \\
\hline
\end{tabular}
\end{center}

\textbf{规则}：看微分式右边的变量——那就是该热力学势的\textbf{独立变量}（自然变量）。

\textbf{例子——熵 $S$ 到底依赖于谁}：
\begin{itemize}[nosep]
  \item 从 $U(S,V,N)$ 看：$S$ 是 $U$ 的\textbf{自变量}，$S = S(U,V,N)$；
  \item 从 $F(T,V,N)$ 看：$S = -(\partial F/\partial T)_{V,N}$，所以 $S = S(T,V,N)$（封闭系统，固定 $N$）；
  \item 从 $\Omega(T,V,\mu)$ 看：$S = -(\partial \Omega/\partial T)_{V,\mu}$，所以 $S = S(T,V,\mu)$（开放系统，固定 $\mu$）。
\end{itemize}

\textbf{关键认知}：$S = S(T,V,N)$ 和 $S = S(T,V,\mu)$ \textbf{并不矛盾}——它们描述的是\textbf{不同约束条件下}的同一个物理量。封闭系统中 $N$ 固定，开放系统中 $\mu$ 固定。就像力学中 $F = ma$ 和 $F = m\ddot{x}$ 并不矛盾，只是用了不同的变量描述同一规律。

\vspace{0.5em}
\textbf{化学势 $\mu$ 的多种定义}：

化学势是粒子数 $N$ 的\textbf{共轭变量}，从不同热力学势出发有不同表达式：
\[
  \mu = \left(\frac{\partial U}{\partial N}\right)_{S,V}
    = \left(\frac{\partial F}{\partial N}\right)_{T,V}
    = \left(\frac{\partial G}{\partial N}\right)_{T,p} .
\]

\textbf{为什么常用 $F$ 来定义}：因为在实际计算中，$F(T,V,N)$ 比 $U(S,V,N)$ 更方便——$T$ 和 $V$ 是实验上容易控制的变量，而 $S$ 不容易直接测量。所以化学势的\textbf{实用定义}通常是 $\mu = (\partial F/\partial N)_{T,V}$。

\vspace{0.5em}
\textbf{热容是复合偏导吗}：
\[
  C_V = \left(\frac{\partial U}{\partial T}\right)_V .
\]
但 $U$ 的自然变量是 $(S,V)$，不是 $(T,V)$。所以 $C_V$ \textbf{不能}直接从 $\mathrm{d}U$ 读出，必须通过链式法则：
\[
  C_V = \left(\frac{\partial U}{\partial S}\right)_V\left(\frac{\partial S}{\partial T}\right)_V
      = T\left(\frac{\partial S}{\partial T}\right)_V .
\]
同理 $C_p = T(\partial S/\partial T)_p$。

\textbf{结论}：热容的定义确实涉及"跨热力学势的复合偏导"。考试中如果题目给出 $U(T,V)$ 的显式表达式，可以直接对 $T$ 求偏导；但如果只给了 $U(S,V)$ 的关系，就必须用链式法则转换。

\vspace{0.5em}
\textbf{考场建议}：拿到热力学题，\textbf{先写出对应热力学势的全微分}，明确谁是独立变量，再决定求导路径。如果需要求 $\partial U/\partial T$ 但 $U$ 的自然变量里没有 $T$，就用链式法则——这是避免漏项的根本方法。
\end{insightbox}
